% MICA — Literature Review (formal academic version)
% Times New Roman 12pt via newtxtext/newtxmath. Compile: tectonic literature_review.tex
\documentclass[12pt]{article}
\usepackage[margin=1in]{geometry}
\usepackage{newtxtext,newtxmath}
\usepackage[hidelinks]{hyperref}
\usepackage{setspace}
\setstretch{1.15}

\title{Literature Review:\\ Intention Recognition and Calibrated Assistance from Raw Observation}
\author{Joosung Park\\ \small Multi-Campus REU 2026 $\cdot$ Penn State University $\cdot$ Computer Science}
\date{July 2026}

\begin{document}
\maketitle

\noindent
This review situates MICA (Minecraft Intention-Aware Collaborative Agent) within
five research threads: assistance games, goal and intent recognition, Minecraft
foundation models, structured action prediction, and confidence calibration.
Although each thread supplies an essential ingredient, none combines explicit
calibrated belief tracking with perception-based evidence under raw observation;
consequently, that combination defines the gap this project addresses.

\section{Assistance Under Hidden Goals}

The formal foundation of this work is the assistance game. Hadfield-Menell et
al.~\cite{cirl} formalized value alignment as Cooperative Inverse Reinforcement
Learning (CIRL), a two-player cooperative game in which both agents are rewarded
according to the human's objective, although only the human observes that
objective. Because the robot must therefore infer the goal from behavior while
simultaneously assisting, the authors prove that isolated optimality is
suboptimal in this setting; thus, an effective assistant must manage its
uncertainty rather than act on a point estimate. However, CIRL specifies the
problem rather than a system.

Puig et al.~\cite{wah} operationalized this framework as the Watch-And-Help
benchmark, in which an agent first observes a human demonstration (Watch) and
subsequently assists with the same goal in a novel environment (Help). Two
findings from this benchmark directly motivate the present work. First, the cost
of incorrect inference is negative rather than merely zero: a helper acting on a
wrong goal repeatedly undid the human's work, and goal conflicts arose in 40\%
of episodes. Second, the authors explicitly identify online goal inference---
updating the estimate during assistance rather than fixing it after a single
demonstration---as the principal open problem. Accordingly, NOPA~\cite{nopa}
extended this line by maintaining a set of goal hypotheses that are updated
online as the human acts, and by modulating assistance according to inference
uncertainty. Most recently, AssistanceZero~\cite{assistancezero} scaled
assistance games to collaborative Minecraft building---the same task family as
MICA---using a Monte Carlo tree search policy whose goal inference is amortized
into a neural prediction head.

Collectively, this thread establishes the problem structure, the cost of
erroneous inference, and the value of explicit uncertainty. Nevertheless, all of
these systems share one simplification: the assistant reads privileged,
ground-truth environment state, whether as symbolic scene graphs~\cite{wah,nopa}
or as oracle game state~\cite{assistancezero}. Moreover, none of them evaluates
whether the stated confidence of the inference is calibrated, and assistance is
assessed by task outcome alone. In contrast, a physically embodied assistant
receives only noisy, partial perception; therefore, it must additionally
determine when its own inference should not be trusted.

\section{Goal and Intent Recognition}

Goal recognition constitutes an older and broader tradition. Classical plan
recognition~\cite{grdsurvey,vered} infers goals by identifying the candidate
whose optimal plan best explains the observed actions; however, its recognized
weakness is that humans rarely behave optimally. Two responses to this weakness
shaped MICA directly. First, GRAIL~\cite{grail} replaces planning with learning
by training one goal-conditioned policy per candidate goal via imitation
learning from potentially suboptimal demonstrations, so that scoring a partial
trajectory requires only a single forward pass; MICA's learned likelihood heads
follow the same principle, since they learn what pursuing each goal looks like
from recorded sessions rather than assuming it. Second, Zhi-Xuan et
al.~\cite{sips} model the human as a boundedly rational, resource-limited
replanner within an online recursive Bayesian filter, and demonstrate that
modeling the planning process---rather than merely action noise---is what makes
goal posteriors accurate on suboptimal and backtracking trajectories.
Consequently, their latent internal-process variable is the direct ancestor of
MICA's reasoning-mode variable $z$, which distinguishes deliberate from habitual
execution.

In parallel, a robotics line based on partially observable Markov decision
processes~\cite{apomdp,saborio,jain,gaze} established the pattern of a belief
over human intent governing the degree of machine autonomy in shared-control
settings. MICA's escalating authority ladder (observe, suggest, preview, place)
renders this principle explicit as a safety gate.

Finally, CLIPS~\cite{clips} resolves a design question that MICA's evaluation
poses directly, namely where a large language model belongs within intent
inference. In CLIPS, the language model serves exclusively as a likelihood term
inside a Bayesian posterior and never as the decision-maker; notably, their
direct-LLM baseline, in which GPT-4V names the goal outright, was the worst
method evaluated, whereas the LLM embedded in the Bayesian frame performed best.
MICA's four-arm comparison mirrors this contrast, since the LLM-reader arm
(0.292) underperforms the learned Bayesian tracker (0.375) on never-seen
sessions. Furthermore, CLIPS handles the assistant's own actions through a
do-operator, and MICA pins the same rule: the agent's own placed blocks are
never treated as evidence about the human.

\section{Minecraft Foundation Models Read in Reverse}

Minecraft became a viable research platform because two works equipped it with
foundation models. MineDojo~\cite{minedojo} contributed an internet-scale
knowledge base together with MineCLIP, a contrastive video--language model that
scores how well recent video matches a natural-language task description.
Meanwhile, VPT~\cite{vpt} trained a behavioral prior on approximately 70{,}000
hours of YouTube gameplay, pseudo-labeled by an inverse dynamics model. An
ecosystem of acting agents subsequently emerged: STEVE-1~\cite{steve1} follows
text instructions by writing commands into MineCLIP's latent space, multi-agent
benchmarks~\cite{teamcraft,villageragent,mbag} coordinate several builders, and
BEDD~\cite{bedd} benchmarks fuzzy, human-judged tasks.

MICA's relationship to this ecosystem is an inversion. Whereas each of these
models was constructed to make an agent act, MICA executes the direction in
reverse: VPT's trunk serves as a frozen behavior encoder that summarizes what
the human is doing, and where STEVE-1 writes a command into MineCLIP's goal
space to produce behavior, MICA reads from the same space by scoring observed
behavior against each candidate goal expressed as text. No backbone is
fine-tuned; instead, only lightweight task-specific heads are trained on the
project's own recorded sessions, because the dataset of 37 real building
sessions is far too small to train models of this scale without memorization.
The same pattern extends to three dimensions, since Uni3D~\cite{uni3d}, a
pretrained point-cloud encoder, reads the reconstructed structure as a single
shape vector, following the structured-3D regime established by works such as
SpatialLM~\cite{spatiallm}.

\section{Structured Action Prediction}

MICA's anticipation result depends on treating builder actions as a small formal
language of 87 tokens---five action verbs together with coordinate and
block-type arguments---and on training a decoder to predict what follows. Two
works license this architecture. FastSceneScript~\cite{fastscenescript}
demonstrates that structured scenes can be decoded as token sequences using
multi-token prediction with confidence-guided decoding, which is the same
structured-language framing MICA applies to building actions. Moreover, Gloeckle
et al.~\cite{mtp} provide the peer-reviewed basis for the decoder's multi-token
form by showing that predicting $K$ future tokens through parallel heads on a
shared trunk is both viable and beneficial. Their appendix also reports that
attaching a multi-token loss to a next-token-pretrained model at fine-tuning
time failed; consequently, MICA's staged training recipe, which performs
next-token pretraining on real sessions before the multi-token stage, was
designed to respect rather than ignore this negative result.

\section{Confidence Calibration}

The final thread is the source of this project's principal negative finding.
Guo et al.~\cite{guo} demonstrated that modern neural networks are
systematically overconfident, in that their stated probabilities exceed their
actual accuracies, and they standardized the measurement protocol of reliability
diagrams and expected calibration error (ECE). MICA adopted this protocol as a
pre-registered evaluation of its own belief, and the predicted failure was
confirmed: above a stated confidence of 0.4, held-out reliability inverts.
Therefore, the confidence-licensed placement route is disabled by its own
threshold, which constitutes the project's central safety mechanism. Notably,
prior assistance systems~\cite{wah,nopa,assistancezero} never performed this
audit; in contrast, MICA treats calibration as a gate that operational authority
must pass before any autonomous action is permitted.

\section{Positioning}

In summary, each neighboring line of work omits precisely one element that MICA
supplies. The assistance-game lineage~\cite{cirl,wah,nopa,assistancezero}
establishes the problem yet reads oracle state and never audits calibration; the
goal-recognition lineage~\cite{grail,sips,clips} contributes learned
likelihoods, bounded rationality, and Bayesian discipline, yet operates in
symbolic domains without perception; the Minecraft foundation
models~\cite{vpt,minedojo,steve1} supply powerful representations, yet were
built for control rather than interpretation; and calibration
research~\cite{guo} defines the measurement, yet has not been applied to an
assistant's license to act. MICA therefore occupies the intersection: an
explicit, recursive Bayesian belief over goal and reasoning
mode~\cite{sips,clips}, whose likelihoods are learned from
recordings~\cite{grail}, over evidence produced by frozen foundation encoders
run in reverse~\cite{vpt,minedojo,uni3d}, in the same task family as scaled
assistance games~\cite{assistancezero}---together with the discipline none of
these systems carried: the belief's confidence is calibration-audited, and the
agent's authority to act is gated on passing that audit, which it has not yet
passed. The literature thus supplies every ingredient; the contribution is
their combination under robot-like observational conditions, evaluated by
pre-registered criteria, with failures reported alongside successes.

\begin{thebibliography}{99}
\setlength{\itemsep}{2pt}

\bibitem{cirl} D. Hadfield-Menell, A. Dragan, P. Abbeel, and S. Russell,
``Cooperative inverse reinforcement learning,'' in \emph{Advances in Neural
Information Processing Systems (NeurIPS)}, 2016.

\bibitem{wah} X. Puig, T. Shu, S. Li, Z. Wang, Y.-H. Liao, J. B. Tenenbaum,
S. Fidler, and A. Torralba, ``Watch-And-Help: A challenge for social perception
and human-AI collaboration,'' in \emph{International Conference on Learning
Representations (ICLR)}, 2021.

\bibitem{nopa} X. Puig, T. Shu, J. B. Tenenbaum, and A. Torralba, ``NOPA:
Neurally-guided online probabilistic assistance for building socially
intelligent home assistants,'' in \emph{IEEE International Conference on
Robotics and Automation (ICRA)}, 2023.

\bibitem{assistancezero} C. Laidlaw, E. Bronstein, T. Guo, D. Feng,
L. Berglund, J. Svegliato, S. Russell, and A. Dragan, ``AssistanceZero:
Scalably solving assistance games,'' in \emph{International Conference on
Machine Learning (ICML)}, 2025.

\bibitem{grdsurvey} S. Keren, A. Gal, and E. Karpas, ``Goal recognition
design -- survey,'' in \emph{International Joint Conference on Artificial
Intelligence (IJCAI), Survey Track}, 2020.

\bibitem{vered} M. Vered and G. A. Kaminka, ``Heuristic online goal recognition
in continuous domains,'' in \emph{International Joint Conference on Artificial
Intelligence (IJCAI)}, 2017.

\bibitem{grail} O. Elhadad, F. Meneguzzi, and R. Mirsky, ``GRAIL: Goal
recognition alignment through imitation learning,'' in \emph{International
Conference on Autonomous Agents and Multiagent Systems (AAMAS)}, 2026.

\bibitem{sips} T. Zhi-Xuan, J. L. Mann, T. Silver, J. B. Tenenbaum, and
V. K. Mansinghka, ``Online Bayesian goal inference for boundedly rational
planning agents,'' in \emph{Advances in Neural Information Processing Systems
(NeurIPS)}, 2020.

\bibitem{apomdp} C.-P. Lam and S. S. Sastry, ``A POMDP framework for
human-in-the-loop system,'' in \emph{IEEE Conference on Decision and Control
(CDC)}, 2014.

\bibitem{saborio} J. C. Saborio and J. Hertzberg, ``Planning with intention
recognition in partially observable domains,'' arXiv preprint, 2024.

\bibitem{jain} S. Jain and B. Argall, ``Probabilistic human intent recognition
for shared autonomy in assistive robotics,'' \emph{ACM Transactions on
Human-Robot Interaction}, 2019.

\bibitem{gaze} V. Belcamino, M. Takase, M. Kilina, A. Carf\`i, A. Shimada,
S. Shimizu, and F. Mastrogiovanni, ``Gaze-based intention recognition for
human-robot collaboration,'' in \emph{ACM/IEEE International Conference on
Human-Robot Interaction (HRI)}, 2024.

\bibitem{clips} T. Zhi-Xuan, L. Ying, V. Mansinghka, and J. B. Tenenbaum,
``Pragmatic instruction following and goal assistance via cooperative
language-guided inverse planning,'' in \emph{International Conference on
Autonomous Agents and Multiagent Systems (AAMAS)}, 2024.

\bibitem{minedojo} L. Fan, G. Wang, Y. Jiang, A. Mandlekar, Y. Yang, H. Zhu,
A. Tang, D.-A. Huang, Y. Zhu, and A. Anandkumar, ``MineDojo: Building
open-ended embodied agents with internet-scale knowledge,'' in \emph{Advances
in Neural Information Processing Systems (NeurIPS)}, 2022.

\bibitem{vpt} B. Baker, I. Akkaya, P. Zhokhov, J. Huizinga, J. Tang,
A. Ecoffet, B. Houghton, R. Sampedro, and J. Clune, ``Video PreTraining (VPT):
Learning to act by watching unlabeled online videos,'' in \emph{Advances in
Neural Information Processing Systems (NeurIPS)}, 2022.

\bibitem{steve1} S. Lifshitz, K. Paster, H. Chan, J. Ba, and S. McIlraith,
``STEVE-1: A generative model for text-to-behavior in Minecraft,'' in
\emph{Advances in Neural Information Processing Systems (NeurIPS)}, 2023.

\bibitem{teamcraft} Q. Long, Z. Li, R. Gong, Y. N. Wu, D. Terzopoulos, and
X. Gao, ``TeamCraft: A benchmark for multi-modal multi-agent Minecraft,'' in
\emph{Advances in Neural Information Processing Systems (NeurIPS)}, 2024.

\bibitem{villageragent} Y. Dong, X. Zhu, Z. Pan, L. Zhu, and Y. Yang,
``VillagerAgent: A graph-based multi-agent framework for coordinating dependent
tasks in Minecraft,'' in \emph{Annual Meeting of the Association for
Computational Linguistics (ACL)}, 2024.

\bibitem{mbag} J. Lee, ``Minecraft building assist game,'' in \emph{IEEE
International Conference on Big Data}, 2022.

\bibitem{bedd} S. Milani, A. Kanervisto, K. Ramanauskas, S. Schulhoff,
B. Houghton, and R. Shah, ``BEDD: The MineRL BASALT evaluation and
demonstrations dataset,'' in \emph{Advances in Neural Information Processing
Systems (NeurIPS), Datasets and Benchmarks Track}, 2023.

\bibitem{uni3d} J. Zhou, J. Wang, B. Ma, Y.-S. Liu, T. Huang, and X. Wang,
``Uni3D: Exploring unified 3D representation at scale,'' in \emph{International
Conference on Learning Representations (ICLR)}, 2024.

\bibitem{spatiallm} Y. Mao, J. Zhong, C. Fang, J. Zheng, R. Tang, H. Zhu,
P. Tan, and Z. Zhou, ``SpatialLM: Large language models for structured indoor
scene understanding,'' in \emph{IEEE/CVF Conference on Computer Vision and
Pattern Recognition (CVPR)}, 2024.

\bibitem{fastscenescript} R. Yin, X. Shi, O. Bailo, M. Manfredi, and
T. Gevers, ``FastSceneScript: Accelerating structured scene generation via
multi-token prediction,'' arXiv preprint arXiv:2512.05597, 2025.

\bibitem{mtp} F. Gloeckle, B. Y. Idrissi, B. Rozi\`ere, D. Lopez-Paz, and
G. Synnaeve, ``Better \& faster large language models via multi-token
prediction,'' in \emph{International Conference on Machine Learning (ICML)},
2024.

\bibitem{guo} C. Guo, G. Pleiss, Y. Sun, and K. Q. Weinberger, ``On
calibration of modern neural networks,'' in \emph{International Conference on
Machine Learning (ICML)}, 2017.

\end{thebibliography}
\end{document}
